\section{Valores numéricos y alcance}\label{sec:discusion}

\subsection{Valores numéricos}

La Tabla~\ref{tab:numerica} recoge una resolución numérica determinista del sistema de Euler
reescalado: método de Newton en mallas espaciales de $64$, $128$ y $256$ puntos con tolerancia
$3\cdot10^{-12}$, una malla en $x$ de paso $0{,}01$ con nodos adicionales, interpolación por
\emph{splines} naturales y regla de Simpson para $\int\log F$; la derivada $F'(0)$ se obtiene por
diferenciación implícita del sistema discretizado. La diferencia máxima entre las curvas de $256$ y
$128$ puntos, en $1001$ nodos, es $1{,}085\cdot10^{-5}$; el residuo de Newton máximo es
$2{,}607\cdot10^{-12}$; la pendiente implícita y la diferencia centrada con paso $10^{-4}$ en $x=0$,
en la malla de $256$ puntos, difieren en $4{,}63\cdot10^{-9}$. Nobel, Agrawal y Boyd~\cite{NAB}
usan también un método de Newton para el problema de masa uno discretizado.

\begin{table}[ht]
\centering
\begin{tabular}{lrrr}
\toprule
Malla & $F(0)$ & $F'(0)$ & $\int_0^1\log F$\\
\midrule
$64$ & $0{,}08086911055$ & $0{,}5157082194$ & $-0{,}9441617254$\\
$128$ & $0{,}08090292111$ & $0{,}5157547494$ & $-0{,}9440411866$\\
$256$ & $0{,}08091137290$ & $0{,}5157663849$ & $-0{,}9440110586$\\
\midrule
Richardson & $0{,}080914$ & $0{,}515770$ & $-0{,}944001$\\
\bottomrule
\end{tabular}
\caption{Valores numéricos de la función $F$. Las diferencias entre mallas consecutivas decrecen
en un factor cercano a cuatro, compatible con convergencia de segundo orden; la última fila es la
extrapolación de Richardson de las mallas de $128$ y $256$ puntos con ese orden, calculada a partir de
las filas anteriores. Son valores numéricos, sin certificación por intervalos.}
\label{tab:numerica}
\end{table}

Por la Proposición~\ref{prop:intlogF}, $\sopt(1)=\int_0^1\log F-1\approx-1{,}9440$. El valor
$-\sopt(1)\approx1{,}9440$ está dentro del intervalo $\alpha=1{,}942\pm3\cdot10^{-3}$
de~\cite{Simkin}, y el valor extrapolado, $-\sopt(1)\approx1{,}944001$, está dentro del intervalo
$[1{,}944000752;\,1{,}944001082]$ de~\cite{NAB}. La Observación~\ref{obs:simkin} explica la normalización que permite comparar ambas
cantidades. La concordancia de los decimales es una comprobación numérica. Las extrapolaciones con las mallas $(64,128)$ y $(128,256)$ difieren en menos de
$10^{-8}$; el error de la cuadratura en $x$ no está estimado.
\ifx\DOIcodigo\empty\else
Un programa que reproduce estos valores con otro método, una iteración de punto fijo, está
en~\cite{Codigo}; las diferencias con la tabla son menores que $10^{-8}$.
\fi
Por simulación de Monte Carlo en tableros de $1024\times1024$, Liu, Liao y Wang~\cite{LLW} estiman
$\alpha\approx1{,}946\pm0{,}003$ integrando el calor específico desde temperatura infinita.

\subsection{Interpretación en el tablero lleno}

Cuando $x=0$, las vacancias de filas y columnas valen cero. Los factores
$r_{R,0}$ y $r_{C,0}$ describen, respectivamente, los límites de las densidades de vacancia
en filas y en columnas, cada una normalizada a masa uno. Junto con las vacancias diagonales
límite y el factor $b_0$ determinan $\nu_1$ por el
Teorema~\ref{thm:lleno}. La regularidad espacial obtenida es Lipschitz, con una misma
cota para $0\le x\le 1/2$. La analiticidad describe una dependencia distinta,
la de los perfiles respecto del parámetro $x$.

Todos los enunciados de este artículo se refieren al problema continuo. Las constantes de las
Secciones~\ref{sec:cotas} y~\ref{sec:regularidad} son explícitas pero muy conservadoras, y no
tienen utilidad numérica. La Tabla~\ref{tab:numerica} estima $F(0)$ resolviendo el sistema continuo.
La identidad $F(0)=1/b_0$ procede del Teorema~\ref{thm:F}. Interpretar este valor
como una frecuencia en tableros finitos exige además especificar la ley de las
colocaciones y justificar el paso del modelo discreto al continuo.
